%%%%%%%%%%%%%%%%%%%%%%%%%%%%%%%%%%%%%%%%%%%%%%%%%%%%%%%%%%%%%%%%%%%%%%%%%%%
%% Entradas del glosario (recomendación: añadir "g" tras la etiqueta).
%%%%%%%%%%%%%%%%%%%%%%%%%%%%%%%%%%%%%%%%%%%%%%%%%%%%%%%%%%%%%%%%%%%%%%%%%%%

\newglossaryentry{GULg}{
    name={Grupo de Usuarios de Linux},
    description={Asociación de estudiantes con inquietudes en torno a los
                 ordenadores y la informática y, en concreto, con la idea común
                 de utilizar y promover el software libre}
}

\newglossaryentry{UC3Mg}{
    name={Universidad Carlos III de Madrid},
    description={Universidad pública española situada en la Comunidad de Madrid}
}

\newglossaryentry{swlibre}{
    name={Software libre},
    description={Software que respeta la libertad de las personas usuarias para
                 ejecutarlo, copiarlo, estudiarlo, modificarlo y redistribuirlo}
}


%%%%%%%%%%%%%%%%%%%%%%%%%%%%%%%%%%%%%%%%%%%%%%%%%%%%%%%%%%%%%%%%%%%%%%%%%%%
%% Entradas de acrónimos (recomendación: añadir "a" tras la etiqueta).
%%  Un acrónimo puede tener su entrada de glosario (GUL, UC3M) mediante see=;
%%  o no tenerla (SO), para ilustrar ambos casos.
%%%%%%%%%%%%%%%%%%%%%%%%%%%%%%%%%%%%%%%%%%%%%%%%%%%%%%%%%%%%%%%%%%%%%%%%%%%

\newglossaryentry{GULa}{
    type = \acronymtype,
    name = {GUL},
    description = {Grupo de Usuarios de Linux},
    see=[Glosario:]{GULg}
}

\newglossaryentry{SOa}{
    type = \acronymtype,
    name = {SO},
    description = {Sistema Operativo}
}

\newglossaryentry{UC3Ma}{
    type = \acronymtype,
    name = {UC3M},
    description = {Universidad Carlos III de Madrid},
    see=[Glosario:]{UC3Mg}
}
